%!Tex Root = ../main.tex
% ./Packete.tex
% ./Design.tex
% ./Vorbereitung.tex
% ./Aufgabe1.tex
% ./Aufgabe2.tex
% ./Aufgabe3.tex
% ./Aufgabe4.tex
% ./Appendix.tex

% ┌──────────┐
% │ Settings │
% └──────────┘

% https://tex.stackexchange.com/questions/325003/given-enumitem-beamer-incompatibility-how-do-i-adjust-the-indent-of-the-enume
% https://tex.stackexchange.com/questions/455692/beamer-presentation-with-itemize-exceed-text-capacity
% \setlist[itemize]{itemsep=2mm, topsep=2mm, parsep=0mm, partopsep=0mm}

% ------------------------------ Minted Settings ------------------------------

% https://tex.stackexchange.com/questions/584071/configure-minted-style-in-latex-for-code-highlighting
% https://pygments.org/docs/tokens/
% https://pygments.org/docs/styledevelopment/#creating-own-styles
% 'custom':  'custom::CustomStyle' in /usr/lib/python3.10/site-packages/pygments/styles/__init__.py
% https://tex.stackexchange.com/questions/18083/how-to-add-custom-c-keywords-to-be-recognized-by-minted#comment930474_42392
% \usemintedstyle{custom}

% https://tex.stackexchange.com/questions/345976/global-settings-for-minted
% \setminted{fontsize=\tiny,breaklines,highlightcolor=SecondaryColorDimmed,autogobble,escapeinside=||,breakafter={_},breakbefore={(},breakaftersymbolpre={},breakaftersymbolpost={},breakbeforesymbolpre={},breakbeforesymbolpost={},breaksymbolsepleft=2mm,breaksymbolsepright=0mm,breakindent=0mm,breaksymbolindentleft=5mm,breaksymbolindentright=0mm,numbersep=2.5mm}
\setminted{fontsize=\tiny,breaklines,highlightcolor=SecondaryColorDimmed,autogobble,escapeinside=||,breakaftersymbolpre={},breakaftersymbolpost={},breakbeforesymbolpre={},breakbeforesymbolpost={},breaksymbolsepleft=2mm,breaksymbolsepright=0mm,breakindent=0mm,breaksymbolindentleft=5mm,breaksymbolindentright=0mm,numbersep=2.5mm}

\newenvironment{linenums}{
  \setminted{linenos}
}{}

% ┌───────┐
% │ Boxes │
% └───────┘

\DeclareTCBInputListing{\codebox}{ s o m }{listing file={#3},
  enhanced,colframe=PrimaryColor,colback=BoxColor,IfBooleanTF={#1}{colframe=SecondaryColor}{colframe=PrimaryColor},fonttitle=\tiny,#2,listing only,halign title=center,drop fuzzy shadow,arc=1mm,bottom=1mm,top=1mm,left=1mm,right=1mm,boxrule=0.5mm,listing engine=minted}
% , sharpish corners

% % https://tex.stackexchange.com/questions/585582/inside-of-a-newtcbinputlisting-how-can-i-change-the-color-of-the-line-number-as
% % https://www.overleaf.com/learn/latex/Using_colours_in_LaTeX
% https://tex.stackexchange.com/questions/132849/how-can-i-change-the-font-size-of-the-number-in-minted-environment
\renewcommand{\theFancyVerbLine}{\ttfamily\textcolor{white}{\tiny{\arabic{FancyVerbLine}}}}
% \renewcommand{\theFancyVerbLine}{\sffamily \textcolor[rgb]{0.5,0.5,1.0}{\huge \oldstylenums{\arabic{FancyVerbLine}}}}

\newtcbinputlisting{\numberedcodebox}[2][]{
  listing file={#2}, enhanced, colframe=PrimaryColor,colback=BoxColor, fonttitle=\small, #1, listing only, halign title=center,drop fuzzy shadow,arc=1mm,boxrule=0.5mm,listing engine=minted,overlay={\begin{tcbclipinterior}\fill[PrimaryColor] (frame.south west) rectangle ([xshift=4mm]frame.north west);\end{tcbclipinterior}}
}

\DeclareTColorBox{codeframe}{ s o m }{
  enhanced, halign title=center, fonttitle=\tiny, interior style={fill=white}, IfBooleanTF={#1}{frame style={color=SecondaryColor}}{frame style={color=PrimaryColor}}, title={#3}, #2,drop fuzzy shadow,arc=1mm,bottom=1mm,top=1mm,left=1mm,right=1mm,boxrule=0.5mm,listing engine=minted,minted style=colorful}

\newtcolorbox{file}[1][]{enhanced, hbox, notitle, interior style={fill=PrimaryColor}, frame empty, halign=center, fontupper=\color{white}\tiny, #1,drop fuzzy shadow,arc=1mm,bottom=1mm,top=1mm,left=1mm,right=1mm,boxrule=0.5mm,listing engine=minted,minted style=colorful}

\newtcblisting{terminal}{
enhanced,colframe=PrimaryColor,colback=BoxColor,hbox,listing only,halign title=center,minted language=text,drop fuzzy shadow,arc=1mm,bottom=1mm,top=1mm,left=1mm,right=1mm,boxrule=0.5mm,listing engine=minted,minted style=colorful,minted options={}}

% % https://tex.stackexchange.com/questions/593218/nested-inline-math-for-new-command-with-argument
\newcommand{\prompt}{\textcolor{SecondaryColor}{\setBold >\;\ignorespaces}}
% % https://tex.stackexchange.com/questions/593218/nested-inline-math-for-new-command-with-argument
\newcommand{\customprompt}{\textnormal\bfseries\textcolor{SecondaryColor}{S}\textcolor{gray!90!black}{he}\textcolor{SecondaryColorDimmed}{ll}\textcolor{SecondaryColor}{>}\;\ignorespaces}

\DeclareTotalTCBox{\inlinebox}{ s v }
{verbatim,colback=SecondaryColorDimmed,colframe=SecondaryColor}
{\IfBooleanTF{#1}%
{\textcolor{SecondaryColor}{\setBold >\enspace\ignorespaces}#2}%
{#2}}

\DeclareTotalTCBox{\key}{ m }
{verbatim,colback=SecondaryColorDimmed,colframe=SecondaryColor}
{$\mathtt{#1}$}

\newtcolorbox{Sidenote}{enhanced,breakable,drop fuzzy shadow,sharpish corners, notitle,arc=0mm,left=3mm,right=3mm,boxrule=1mm, borderline vertical={1mm}{0pt}{PrimaryColor},title=Sidenote,attach boxed title to top text left={yshift=0mm},
interior style={fill=BoxColor},
frame style={color=BoxColor},
% https://tex.stackexchange.com/questions/459870/paragraph-breaks-inside-tcolorbox, maybe also parbox=false
boxed title style={arc=0mm,skin=enhancedfirst jigsaw,frame empty,colback=PrimaryColor,boxrule=0mm,bottom=-0.4mm},
after title={\hspace{0.2cm}\includegraphics[height=3mm]{./figures/lupe.png}}
}
% before upper=\setlength{\parskip}{1em},
% fonttitle=\bfseries

\numberwithin{equation}{section}

% ┌───────────────┐
% │ Beamer Blocks │
% └───────────────┘

\newcounter{exercise}
\setcounter{exercise}{1}

% https://tex.stackexchange.com/questions/309463/beamer-numerating-theorems-and-color-of-block
\setbeamertemplate{theorems}[numbered]

% https://tex.stackexchange.com/questions/152565/define-a-new-block-environment-in-latex-beamer
\renewenvironment<>{exercise}[1][]{%
  \setbeamercolor{block title}{fg=white,bg=SecondaryColor}%
\begin{block}#2{Exercise \thesection{.}\theexercise\enspace\ignorespaces #1 \anotebook}}{\end{block}\stepcounter{exercise}}

% https://tex.stackexchange.com/questions/152565/define-a-new-block-environment-in-latex-beamer
\renewenvironment<>{exercisenoinc}[1][]{%
  \setbeamercolor{block title}{fg=white,bg=SecondaryColor}%
\begin{block}#2{
Exercise \thesection{.}\theexercise\enspace\ignorespaces #1 \anotebook}}{\end{block}}

\newenvironment<>{requirements}[1][]{%
  \setbeamercolor{block title}{fg=white,bg=SecondaryColor}%
\begin{block}#2{Requirements \thesection{.}\theexercise\enspace\ignorespaces #1 \anotebook}\itshape}{\end{block}\stepcounter{exercise}}

\newenvironment<>{requirementsnoinc}[1][]{%
  \setbeamercolor{block title}{fg=white,bg=SecondaryColor}%
\begin{block}#2{Requirements \thesection{.}\theexercise\enspace\ignorespaces #1 \anotebook}\itshape}{\end{block}}

\newenvironment<>{solution}[1][]{%
  \setbeamercolor{block title}{fg=white,bg=PrimaryColor}%
\begin{block}#2{
\if\matthias0{
Solution
}\else{
Lösungsidee
}\fi
 \thesection{.}\theexercise\enspace\ignorespaces #1 \notebook}\itshape}{\end{block}\stepcounter{exercise}}

\newenvironment<>{solutionnoinc}[1][]{%
  \setbeamercolor{block title}{fg=white,bg=PrimaryColor}%
\begin{block}#2{
\if\matthias0{
Solution
}\else{
Lösungsidee
}\fi
 \thesection{.}\theexercise\enspace\ignorespaces #1 \notebook}\itshape}{\end{block}}

\newcounter{definition}
\setcounter{definition}{1}

\newenvironment<>{definition}[1]{%
  \setbeamercolor{block title}{fg=white,bg=SecondaryColorDimmed}%
\begin{block}#2{Definition \thedefinition\enspace\ignorespaces #1 \anotebook}\itshape}{\end{block}\stepcounter{definition}}

% https://stackoverflow.com/questions/34274267/how-can-i-number-theorems-definitions-etc-etc-consecutively-in-latex-without
\newtheorem{customtheorem}{Custom Theorem}[section]
\newtheorem{subcustomtheorem}{Sub Custom Theorem}[customtheorem]
\newtheorem{ocustomtheorem}[customtheorem]{Custom Theorem}

% ┌──────────────┐
% │ New Commands │
% └──────────────┘

% alternative alert
\newcommand\aalert[1]{\textcolor{PrimaryColor}{#1}}

\newcommand{\aitem}{%
\item[\textcolor{PrimaryColor}{$\blacktriangleright$}]\ignorespaces}

% % https://tex.stackexchange.com/questions/8351/what-do-makeatletter-and-makeatother-do
% \makeatletter
%
% % ignorespace: https://runebook.dev/de/docs/latex/_005cignorespaces-_0026-_005cignorespacesafterend
% % enspace: https://math-linux.com/latex-26/faq/latex-faq/article/latex-horizontal-space-qquad-hspace-thinspace-enspace
% \newcommand{\aitem[1]}[]{%
% \@ifmtarg{#1}{\item[\textcolor{SecondaryColor}{\blacktriangleright}]}%
% {\item[\textcolor{SecondaryColor}{\black}] \colorbold{\textbf{#1:}}}\enspace\ignorespaces}
% % \textbullet
%
% \makeatother

\newcommand{\grayitem}{%
\item[\textcolor{gray!90!black}{$\blacktriangleright$}]\ignorespaces}

% https://tex.stackexchange.com/questions/329990/how-do-i-change-the-color-of-itemize-bullet-specific-and-default
% \renewcommand{\labelitemi}{$\textcolor{PrimaryColor}{\bullet}$}
% \renewcommand{\labelitemii}{$\textcolor{PrimaryColor}{\cdot}$}
% \renewcommand{\labelitemiii}{$\textcolor{PrimaryColor}{\diamond}$}
% \renewcommand{\labelitemiv}{$\textcolor{PrimaryColor}{\ast}$}

\newcommand{\notebook}{\hfill\includegraphics[height=3mm]{./figures/note.png}}
\newcommand{\anotebook}{\hfill\includegraphics[height=3mm]{./figures/anote.png}}

\newcommand{\magnifier}{\hfill\includegraphics[height=3mm]{./figures/lupe.png}}

\newcommand\RemoveMargin[2][3em]{%
\makebox[\linewidth][c]{%
  \begin{minipage}{\dimexpr\textwidth+#1\relax}
  \raggedright#2
  \end{minipage}%
  }%
}

% https://tex.stackexchange.com/questions/188379/theorem-numbering-in-beamer
% Nummerierung der Theorem und Berücksichtigung der Sections
\renewcommand\thetheorem{\arabic{section}.\arabic{theorem}}

\newcommand{\ma}[1]{$\mathcal{#1}$}
\renewcommand{\tt}[1]{{\small\texttt{#1}}}
% \renewcommand{\pr}[1]{{\textnormal{\textsc{#1}}}}

% ┌──────────────────┐
% │ New Environments │
% └──────────────────┘


\newenvironmentx{transformation}[3][1=0.3, 2=0.2, 3=0.4]{
  \newcommand{\compilesto}{%
    \end{column}\begin{column}{#2\linewidth}\centering$\Rightarrow$\end{column}\begin{column}{#3\linewidth}\centering}
  \newcommand{\arrowx}[1]{%
    \end{column}\begin{column}{#2\linewidth}\centering$\xRightarrow{##1}$\end{column}\begin{column}{#3\linewidth}\centering}
  \newcommand{\arrowxx}[2]{%
    \end{column}\begin{column}{#2\linewidth}\centering$\xRightarrow[##2]{##1}$\end{column}\begin{column}{#3\linewidth}\centering}
  \begin{columns}\begin{column}{#1\linewidth}\centering
  }{%
  \end{column}\end{columns}%
}

\newcounter{algorithm}
\setcounter{algorithm}{0}
\newtcbtheorem[use counter=algorithm]{algorithm}{\color{SecondaryColor}Algorithm}{pseudo/ruled}{alg}
